%------------------------------------------------------------------
%------------------------------------------------------------------
%------------------------------------------------------------------

\chapter{Command-line arguments}

\section{Local serializable enums}

MMVII does not require every enumerated type to live in
\texttt{MMVII\_enums.h}.  A command can define its own enum locally and still
benefit from the usual string conversion, command-line parsing and help display.

The enum declaration must still follow the usual MMVII convention: the last
enumerator is \texttt{eNbVals}.  The string mapping is then provided once, in one
translation unit, through \texttt{cE2Str<Enum>::mE2S},
\texttt{MACRO\_DECLARE\_STRIO\_ENUM} and
\texttt{MACRO\_INSTANTIATE\_STRIO\_ENUM}.  After that, the enum behaves like a
regular MMVII enum for textual conversion.

The following example keeps the enum local to one \texttt{.cpp} file:

\begin{lstlisting}[language=C++]
#include "MMVII_Tpl_ElemStrToVal.h"

enum class eDispExif
{
    eSizeType=0,
    eMainExif,
    eAllExif,
    eRawExif,
    eAllGDALInfo,
    eNbVals
};

template<> cE2Str<eDispExif>::tMapE2Str cE2Str<eDispExif>::mE2S
{
    {eDispExif::eSizeType,"SizeType"},
    {eDispExif::eMainExif,"MainExif"},
    {eDispExif::eAllExif,"AllExif"},
    {eDispExif::eRawExif,"RawExif"},
    {eDispExif::eAllGDALInfo,"AllGDALInfo"}
};

MACRO_INSTANTIATE_STRIO_ENUM(eDispExif,"DispExif")
\end{lstlisting}

The same pattern can be split between a header and a source file if the enum is
used in several compilation units.  The header then carries nothing but the enum
declaration itself, while the string mapping and the instantiation stay in exactly
one \texttt{.cpp} file:

\begin{lstlisting}[language=C++]
// toto.h
enum class eEpipFrm
{
    eIntersect,
    eUnion,
    eImg_1,
    eImg_2,
    eNbVals
};
\end{lstlisting}

\begin{lstlisting}[language=C++]
// toto.cpp
#include "toto.h"

template<> cE2Str<eEpipFrm>::tMapE2Str cE2Str<eEpipFrm>::mE2S
{
    {eEpipFrm::eIntersect,"Intersect"},
    {eEpipFrm::eUnion,"Union"},
    {eEpipFrm::eImg_1,"Img_1"},
    {eEpipFrm::eImg_2,"Img_2"}
};

MACRO_INSTANTIATE_STRIO_ENUM(eEpipFrm,"EpipFrame")
\end{lstlisting}

Once instantiated, the enum can be used directly in a command argument.  MMVII
automatically adds the list of allowed values for enum-typed arguments, so there
is no need to add \texttt{AC\_ListVal} manually in \texttt{Arg2007} or
\texttt{AOpt2007}.  If the enum has not been instantiated, the first use in a
command-line argument or in a string conversion context triggers a compilation
error, which is the intended failure mode.

%------------------------------------------------------------------

\section{Structured command-line arguments}

\subsection{Purpose}

A command-line parameter usually initializes one scalar value or one
\texttt{std::vector}.  Some commands, however, need a group of related values that
should be parsed, documented and completed as a single parameter.  MMVII supports
such parameters through \texttt{ARG2007\_STRUCT\_FIELDS}.

This mechanism is intended for command-line parsing.  It is independent from the
archive serialization mechanism described in the chapter on serialization.

A structured argument provides:

\begin{itemize}
  \item conversion between the structure and its bracketed textual representation;
  \item field metadata for the integrated help and bash completion;
  \item an \texttt{eTA2007} semantic list for each field;
  \item recursive support for structures and vectors of structures;
  \item optional fields with default values.
\end{itemize}

The declarations are available through
\texttt{MMVII\_Tpl\_ElemStrToVal.h}.

%------------------------------------------------------------------

\subsection{Declaring a structured argument}

The following example defines one argument containing an orientation name, an
operating mode, a scale and an output image:

\begin{lstlisting}[language=C++]
#include "MMVII_Tpl_ElemStrToVal.h"

enum class eExportMode
{
    Fast,
    Accurate,
    eNbVals
};

MACRO_INSTANTIATE_STRIO_ENUM(eExportMode,"ExportMode")

struct cExportOptions
{
    std::string Orientation;
    eExportMode Mode;
    double      Scale=1.0;
    int         Offset;
    std::string Output = "Result.tif";

    ARG2007_STRUCT_FIELDS
    (
        Orientation,FieldSem(eTA2007::Orient),
        Mode,
        Scale,FieldSem
               (
                   {
                       eTA2007::HDV,
                       {eTA2007::AllowedValues,"[0.1,1.0,10.0]"},
		       {eTA2007::AddCom,"Description for field Scale"}
                   }
               ),
	Offset,
        Output,FieldSem({eTA2007::HDV,eTA2007::FileImage})
    )
};
\end{lstlisting}

Each ordinary expression in \texttt{ARG2007\_STRUCT\_FIELDS} designates one
field.  Fields are parsed in declaration order.  A \texttt{FieldSem} expression
applies to the field immediately preceding it and does not consume a value from
the command line. The \texttt{FieldSem} expression can be omitted if the field 
has no special semantics.

The macro generates the textual conversion and metadata methods required by
\texttt{Arg2007} and \texttt{AOpt2007}.  The generated textual representation of
the example is:

\begin{verbatim}
[Orientation,Mode,Scale,Offset,Output]
\end{verbatim}

Consequently, this mechanism must only be used once in a given structure, and all
values that form the command-line representation must be listed in the macro.

%------------------------------------------------------------------

\subsection{Field semantics}

\texttt{FieldSem} accepts the same \texttt{eTA2007} attributes that are used for
top-level command arguments.  The supported forms include:

\begin{lstlisting}[language=C++]
FieldSem(eTA2007::HDV)
FieldSem({eTA2007::HDV})
FieldSem({eTA2007::AddCom,"A description"})
FieldSem
(
    {
        eTA2007::HDV,
        {eTA2007::AllowedValues,"[A,B,C]"}
    }
)
\end{lstlisting}

The first two forms contain one semantic.  The third form contains one semantic
with its auxiliary value.  The last form attaches several semantics to one field.

For an enum field, MMVII automatically adds the list of enum values to the field
metadata.  It is therefore unnecessary to duplicate that list with
\texttt{eTA2007::AllowedValues}.

One semantic is explicitly refused on a field: \texttt{eTA2007::CanRepeat}.  It
describes an argument given several times on the command line, which has no meaning
for a field written once, inside its structure.  The refusal is an
\texttt{MMVII\_INTERNAL\_ASSERT\_always} raised while the specification is built, so
\texttt{MMVII GenArgsSpec}, which builds the specification of every command, is the
exhaustive check.

%------------------------------------------------------------------

\subsection{Optional fields}

A field carrying \texttt{eTA2007::HDV} has a default value and may be omitted from
the textual representation.  Its current C++ value is preserved when it is
omitted.  The structure should therefore initialize every field carrying
\texttt{HDV} before parsing.

For the \texttt{cExportOptions} example, all of these values are valid:

\begin{verbatim}
[Ori-Initial,Fast,2.0,1,Export.tif]
[Ori-Initial,Fast,,1,Export.tif]
[Ori-Initial,Fast,,1,]
[Ori-Initial,Fast,,1]
\end{verbatim}


%------------------------------------------------------------------

\subsection{Using the structure in a command}

A structured value is registered exactly like another supported argument type:

\begin{lstlisting}[language=C++]
class cAppliExport : public cMMVII_Appli
{
    // ...
    cExportOptions mOptions;
};

cCollecSpecArg2007 & cAppliExport::ArgObl
                     (cCollecSpecArg2007 & anArgObl)
{
    return anArgObl
           << Arg2007(mOptions,"Export options");
}
\end{lstlisting}

It can also be a named optional argument:

\begin{lstlisting}[language=C++]
cCollecSpecArg2007 & cAppliExport::ArgOpt
                     (cCollecSpecArg2007 & anArgOpt)
{
    return anArgOpt
           << AOpt2007(mOptions,"Options","Export options");
}
\end{lstlisting}

%------------------------------------------------------------------

\subsection{Nested structures and vectors}

A field can itself be a structured argument:

\begin{lstlisting}[language=C++]
struct cExportJob
{
    cExportOptions Export;
    bool           Verbose = false;

    ARG2007_STRUCT_FIELDS
    (
        Export,
        Verbose,FieldSem(eTA2007::HDV)
    )
};
\end{lstlisting}

A corresponding value can be written as:

\begin{verbatim}
[[Ori-Initial,Fast,,1],true]
[[Ori-Initial,Fast,,1]]
\end{verbatim}

The same conversion is applied recursively.  A
\texttt{std::vector<cExportOptions>} uses one additional bracket level:

\begin{verbatim}
[[Ori-1,Fast,,2],[Ori-2,Accurate,2.0,4,Second.tif]]
\end{verbatim}

The element type of a vector of structures must be default-constructible because
MMVII creates a value-initialized element to obtain its field metadata.

A field may also be a \texttt{std::vector} of scalars, which adds its own bracket
level in the same way:

\begin{lstlisting}[language=C++]
struct cSmoothJob
{
    std::string         Mode;
    std::vector<double> Weights;
    int                 NbIter = 1;

    ARG2007_STRUCT_FIELDS
    (
        Mode,
        Weights,
        NbIter,FieldSem(eTA2007::HDV)
    )
};
\end{lstlisting}

\begin{verbatim}
[Gauss,[0.25,0.5,0.25],4]
[Gauss,[1.0]]
\end{verbatim}

Prefer such a field to a \texttt{std::string} holding a bracketed encoding split again
at run time.  The generated metadata then carries the real type, so the help and the
completion show \texttt{[double,...]} instead of \texttt{string}, and a value written
without its brackets becomes a parse error instead of silently shifting into the next
field.

%------------------------------------------------------------------

\subsection{A real-world example}

The bundle-adjustment command \texttt{OriBundleAdj}
(\texttt{src/BundleAdjustment/cAppliBundAdj.cpp}) uses this mechanism for most of
its optional constraints. Its \texttt{cGCP3D} structure, for instance:

\begin{lstlisting}[language=C++]
struct cGCP3D
{
    std::string GCPDir;
    double      Sigma;
    std::string OutDir = "";
    double      ExportSigma = 0;

    ARG2007_STRUCT_FIELDS
    (
        GCPDir,FieldSem({{eTA2007::AddCom,"World coord GCP dir"},eTA2007::ObjCoordWorld}),
        Sigma,FieldSem({{eTA2007::AddCom,"Sigma factor SG=0 fix, SG<0 schurr elim, SG>0"},eTA2007::ObjCoordWorld}),
        OutDir,FieldSem({eTA2007::HDV,{eTA2007::AddCom,"Optional output dir"}}),
        ExportSigma,FieldSem({eTA2007::HDV,{eTA2007::AddCom,"Optional compensated sigma"}})
    )
};
\end{lstlisting}

is registered as a repeatable vector of structures:

\begin{lstlisting}[language=C++]
std::vector<cGCP3D> mGCP3D;
...
<< AOpt2007(mGCP3D,"GCP3D","GCP ground",{eTA2007::CanRepeat})
\end{lstlisting}

Combining a vector of structures with \texttt{eTA2007::CanRepeat} lets the same
command-line key be given several times, each occurrence accumulating into the
vector instead of replacing it:

\begin{verbatim}
GCP3D=[GcpDirA,1.0] GCP3D=[GcpDirB,0.5,OutDirB]
\end{verbatim}

Each occurrence above carries a single element and is written with one bracket level
less than a complete vector value.  That tolerance belongs to
\texttt{eTA2007::CanRepeat} alone, and only for a top-level argument: an argument
written once is given its whole value, so a vector argument without that semantic
must always be fully bracketed.  With \texttt{CanRepeat} the fully bracketed form
stays valid as well, and may carry several elements at once:

\begin{verbatim}
GCP3D=[[GcpDirA,1.0],[GcpDirB,0.5,OutDirB]]
\end{verbatim}

The bash completion applies the very same rule, and offers the shortened form only
for an argument whose specification carries \texttt{CanRepeat}.

%------------------------------------------------------------------

\subsection{Generated help and specification}

The integrated help displays the structured type and a separate field list.  A
question mark identifies each field carrying \texttt{HDV}.  For example:

\begin{verbatim}
  * [string,ExportMode,double?,string?] :: Export options
    - Fields=[Orientation,Mode,Scale?,Offset,Output?]
\end{verbatim}

\texttt{GenArgsSpec} writes the same structure recursively in the \texttt{fields}
property of its JSON output.  Each field contains its name, type and semantics.
Default values and enum values are included when applicable.  This information is
also consumed by the Bash completion script.
